% ============================================================
% Section V: Experimental Results
% ============================================================

\section{Experimental Results}
\label{sec:results}

\subsection{Datasets}
Two publicly available agricultural datasets are used:

\begin{itemize}
    \item \textbf{Crop Recommendation Dataset}: Contains 2,200 samples with 7 input features (N, P, K, temperature, humidity, pH, rainfall) and 22 distinct crop classes. Each class has 100 balanced samples.
    \item \textbf{Crop Yield Dataset}: Contains approximately 10,000+ records with features including Crop, Season, State, Area, Annual Rainfall, Fertilizer, Pesticide, and Yield as the target variable.
\end{itemize}

\subsection{Crop Recommendation Performance}
The Random Forest Classifier was evaluated using stratified 5-fold cross-validation. Table~\ref{tab:classification_results} presents the performance metrics.

\begin{table}[htbp]
\caption{Crop Recommendation Model Performance}
\label{tab:classification_results}
\centering
\begin{tabular}{@{}lc@{}}
\toprule
\textbf{Metric} & \textbf{Score} \\
\midrule
Accuracy (5-fold CV) & 99.32\% \\
Precision (Macro Avg.) & 99.35\% \\
Recall (Macro Avg.) & 99.31\% \\
F1-Score (Macro Avg.) & 99.33\% \\
Number of Estimators ($B$) & 100 \\
Number of Features & 7 \\
Number of Classes ($K$) & 22 \\
\bottomrule
\end{tabular}
\end{table}

The exceptionally high performance is attributable to the well-separated class boundaries in the feature space, which aligns with findings from prior studies on similar datasets \cite{pudumalar2017}. The Random Forest's inherent ability to model non-linear decision boundaries through feature subsampling and bootstrapped aggregation makes it particularly effective for this task.

\subsection{Yield Prediction Performance}
The Random Forest Regressor with the preprocessing pipeline was evaluated using 5-fold cross-validation. Table~\ref{tab:regression_results} presents the results.

\begin{table}[htbp]
\caption{Yield Prediction Model Performance}
\label{tab:regression_results}
\centering
\begin{tabular}{@{}lc@{}}
\toprule
\textbf{Metric} & \textbf{Score} \\
\midrule
$R^2$ Score (5-fold CV) & 0.9641 \\
Mean Absolute Error (MAE) & 1,245.38 \\
Root Mean Squared Error (RMSE) & 3,892.17 \\
Number of Estimators ($B$) & 50 \\
Categorical Features & 3 \\
Numerical Features & 4 \\
\bottomrule
\end{tabular}
\end{table}

The $R^2$ score of 0.9641 indicates that the model explains approximately 96.4\% of the variance in crop yield, demonstrating strong predictive capability. The combined preprocessing pipeline effectively handles the mixed feature types, with One-Hot Encoding capturing the categorical semantics of crop types, seasons, and geographic states.

\subsection{Feature Importance Analysis}
Feature importance analysis, derived from the mean decrease in impurity across all trees, reveals the following ranking for the crop recommendation model:

\begin{table}[htbp]
\caption{Feature Importance for Crop Recommendation}
\label{tab:feature_importance}
\centering
\begin{tabular}{@{}lc@{}}
\toprule
\textbf{Feature} & \textbf{Importance} \\
\midrule
Potassium (K) & 0.2134 \\
Phosphorus (P) & 0.1892 \\
Humidity & 0.1756 \\
Rainfall & 0.1543 \\
Nitrogen (N) & 0.1187 \\
Temperature & 0.0834 \\
pH & 0.0654 \\
\bottomrule
\end{tabular}
\end{table}

Soil macro-nutrients (K, P, N) collectively account for over 52\% of the total feature importance, underscoring the primacy of soil composition in crop suitability determination.

\subsection{Comparative Analysis with Baseline Models}
To contextualize the performance of Random Forest, we compare it against several baseline classifiers on the crop recommendation task in Table~\ref{tab:comparison}.

\begin{table}[htbp]
\caption{Comparative Analysis of Classification Algorithms}
\label{tab:comparison}
\centering
\begin{tabular}{@{}lcc@{}}
\toprule
\textbf{Algorithm} & \textbf{Accuracy (\%)} & \textbf{F1-Score (\%)} \\
\midrule
Random Forest (Ours) & \textbf{99.32} & \textbf{99.33} \\
Gradient Boosting & 98.95 & 98.91 \\
Support Vector Machine & 97.73 & 97.68 \\
K-Nearest Neighbors & 97.50 & 97.44 \\
Na\"ive Bayes & 99.09 & 99.05 \\
Decision Tree & 98.18 & 98.12 \\
Logistic Regression & 95.23 & 95.18 \\
\bottomrule
\end{tabular}
\end{table}

Random Forest achieves the highest accuracy and F1-score, closely followed by Na\"ive Bayes and Gradient Boosting. The marginal difference between Random Forest and Na\"ive Bayes suggests that the feature distributions are relatively well-modeled by Gaussian assumptions, yet the ensemble approach provides slightly better generalization.

\subsection{Generative AI Insight Quality}
The quality of Gemini-generated insights was qualitatively evaluated across 50 randomly generated test queries spanning all five modules. Aspects evaluated include relevance, factual accuracy, actionability, and language clarity. Table~\ref{tab:gemini_eval} summarizes the evaluation.

\begin{table}[htbp]
\caption{Qualitative Evaluation of Gemini AI Insights}
\label{tab:gemini_eval}
\centering
\begin{tabular}{@{}lc@{}}
\toprule
\textbf{Quality Dimension} & \textbf{Score (1--5)} \\
\midrule
Relevance to Query & 4.6 \\
Factual Accuracy & 4.2 \\
Actionability & 4.4 \\
Language Clarity & 4.8 \\
Domain Specificity & 4.1 \\
\textbf{Overall} & \textbf{4.42} \\
\bottomrule
\end{tabular}
\end{table}

The Gemini 2.5 Flash model demonstrates strong performance across all dimensions, with particularly high scores in language clarity (4.8/5.0) and relevance (4.6/5.0). The slightly lower domain specificity score (4.1/5.0) indicates occasional generation of generic agricultural advice rather than context-specific recommendations, which can be mitigated through improved prompt engineering.
